\subsection{附注与进一步结果}
\label{sec:ch7-notes-further-results}

\subsubsection{参考文献}
\label{subsec:ch7-references}

$A_p$ 条件最早以稍有不同的形式见于 M. Rosenblum 的论文 \MBUnresolvedReference{citation}{[Summability of Fourier series in $L^p(\mu)$, Trans. Amer. Math. Soc. 105 (1962), 32--42]}. 当 $n=1$ 时, $A_p$ 的刻画归功于 B. Muckenhoupt, 见 \MBUnresolvedReference{citation}{[Weighted norm inequalities for the Hardy maximal function, Trans. Amer. Math. Soc. 165 (1972), 207--226]}. B. Muckenhoupt, R. Hunt 和 R. Wheeden 在 \MBUnresolvedReference{citation}{[Weighted norm inequalities for the conjugate function and the Hilbert transform, Trans. Amer. Math. Soc. 176 (1973), 227--251]} 中证明: Hilbert 变换在 $L^p(w)$ 上有界当且仅当 $w\in A_p$. 当 $p=2$ 时, H. Helson 和 G. Szegő 给出了另一种刻画, 详见下文第~\ref{subsec:ch7-helson-szego}~节. R. Coifman 和 C. Fefferman 在 \MBUnresolvedReference{citation}{[Weighted norm inequalities for maximal functions and singular integrals, Studia Math. 51 (1974), 241--250]} 中把这些结果推广到高维, 详见下文第~\ref{subsec:ch7-coifman-fefferman}~节. 本章的证明归功于 Journé \MBUnresolvedReference{citation}{[8]}.

\hypertarget{chapter-07-page-160}{}
在同一篇论文中, Coifman 和 Fefferman 还证明了 $A_p$ 权满足关键的反向 Hölder 不等式. 该不等式最早见于 F. W. Gehring 的论文 \MBUnresolvedReference{citation}{[The $L^p$-integrability of the partial derivatives of a quasiconformal mapping, Acta Math. 130 (1973), 3--4, 265--277]}. 极大函数的强型 $(p,p)$ 范数不等式也可不用反向 Hölder 不等式证明, 见 M. Christ 和 R. Fefferman 的 \MBUnresolvedReference{citation}{[A note on weighted norm inequalities for the Hardy--Littlewood maximal operator, Proc. Amer. Math. Soc. 87 (1983), 447--448]}, 以及下文第~\ref{subsec:ch7-two-weights}~节. 定理~\ref{thm:ch7-a1-characterization} 对 $A_1$ 权的刻画归功于 R. Coifman 和 R. Rochberg, 见 \MBUnresolvedReference{citation}{[Another characterization of BMO, Proc. Amer. Math. Soc. 79 (1980), 249--254]}.

外推定理由 J. L. Rubio de Francia 在 \MBUnresolvedReference{citation}{[Factorization theory and $A_p$ weights, Amer. J. Math. 106 (1984), 533--547]} 中证明; 他的证明利用了加权范数不等式与向量值不等式的联系. J. García-Cuerva 在 \MBUnresolvedReference{citation}{[An extrapolation theorem in the theory of $A_p$-weights, Proc. Amer. Math. Soc. 87 (1983), 422--426]} 中给出了不使用向量值不等式的直接证明. 本章的证明更为简单, 至少在组织方式上似乎是新的, 但它遵循标准方法. García-Cuerva 和 Rubio de Francia 的著作 \MBUnresolvedReference{citation}{[6]} 对加权范数不等式的相关结果有出色论述. 另见 B. Muckenhoupt 的综述 \MBUnresolvedReference{citation}{[Weighted norm inequalities for classical operators, Harmonic Analysis in Euclidean Spaces, G. Weiss and S. Wainger, eds., vol. 1, pp. 69--84, Proc. Sympos. Pure Math. 35, Amer. Math. Soc., Providence, 1979]} 以及 E. M. Dynkin 和 B. P. Osilenker 的综述 \MBUnresolvedReference{citation}{[Weighted estimates for singular integrals and their applications, J. Soviet Math. 30 (1985), 2094--2154; translated from Itogi Nauki i Tekhniki, Akad. Nauk SSSR, Moscow, 1983]}.

\subsubsection{Helson--Szegő 条件}
\label{subsec:ch7-helson-szego}

权关于 Hilbert 变换的第一个刻画归功于 H. Helson 和 G. Szegő, 见 \MBUnresolvedReference{citation}{[A problem in prediction theory, Ann. Math. Pura Appl. 51 (1960), 107--138]}: Hilbert 变换在 $L^2(w)$ 上有界当且仅当
\[
 \log w=u+Hv,
\]
其中 $u,v\in L^\infty$ 且 $\|v\|_\infty<\pi/2$. 他们的证明使用复分析. 由该条件推出 $A_2$ 条件是直接的, 但尚无反向蕴含的直接证明.

\hypertarget{chapter-07-page-161}{}
目前最多只能证明上述分解在 $\|v\|_\infty<\pi$ 时成立. 这首先由 R. Coifman, P. Jones 和 J. L. Rubio de Francia 证明, 见 \MBUnresolvedReference{citation}{[Constructive decomposition of BMO functions and factorization of $A_p$ weights, Proc. Amer. Math. Soc. 87 (1983), 675--676]}. J. Garnett 和 P. Jones 给出了这一结果的高维类似物, 见 \MBUnresolvedReference{citation}{[The distance in BMO to $L^\infty$, Ann. of Math. 108 (1978), 373--393]}.

M. Cotlar 和 C. Sadosky 把 Helson--Szegő 定理的条件推广到 $L^p(w)$, 见 \MBUnresolvedReference{citation}{[On some $L^p$ versions of the Helson--Szegő theorem, Conference on Harmonic Analysis in Honor of Antoni Zygmund, W. Beckner et al., eds., vol. 1, pp. 306--317, Wadsworth, Belmont, 1983]}.

\subsubsection{$A_\infty$ 条件}
\label{subsec:ch7-ainfinity-condition}

如前所述, 不等式 \eqref{eq:ch7-ainfinity-condition} 称为 $A_\infty$ 条件, 因为
\[
 A_\infty=\bigcup_{p<\infty}A_p.
\]
这一事实由 B. Muckenhoupt 在 \MBUnresolvedReference{citation}{[The equivalence of two conditions for weight functions, Studia Math. 49 (1974), 101--106]} 中以及 Coifman 和 Fefferman 在上引论文中分别独立证明. 已知 $A_p\subset A_\infty$, 所以只需证明 $w\in A_\infty$ 蕴含存在 $p$ 使 $w\in A_p$. 证明的关键是: $w^{-1}$ 关于测度 $w\,dx$ 满足反向 Hölder 不等式. 即存在正数 $\varepsilon$ 和 $C$, 使每个立方体 $Q$ 都满足
\begin{equation}
\label{eq:ch7-dual-measure-reverse-holder}
 \left(\frac1{w(Q)}\int_Qw^{-1-\varepsilon}w\right)^{1/(1+\varepsilon)}
 \le\frac C{w(Q)}\int_Qw^{-1}w.
\end{equation}
这等价于
\[
 \frac1{|Q|}\int_Qw^{-\varepsilon}
 \le C\left(\frac{|Q|}{w(Q)}\right)^\varepsilon,
\]
即取 $\varepsilon=p'-1$ 时的 $A_p$ 条件.

证明 \eqref{eq:ch7-dual-measure-reverse-holder} 的方法与定理~\ref{thm:ch7-reverse-holder} 的证明类似. 从 $A_\infty$ 条件出发, 可证明引理~\ref{lem:ch7-small-set-weight-decay} 的测度互换版本: 存在 $\alpha'<1$, $\beta'<1$, 使 $S\subset Q$ 且 $w(S)\le\alpha'w(Q)$ 时有 $|S|\le\beta'|Q|$. 注意, 定理~\ref{thm:ch7-reverse-holder} 的证明只用了引理中某一对 $\alpha,\beta$ 的存在性.

完成证明还要用到 $w\in A_\infty$ 时 $w$ 是倍增测度这一事实: 存在 $C>0$, 使任意立方体 $Q$ 都满足 $w(2Q)\le Cw(Q)$. 该性质足以使我们关于测度 $w\,dx$ 在高度 $\lambda$ 对函数作 Calderón--Zygmund 分解.

\hypertarget{chapter-07-page-162}{}
由此得到两两不交的立方体族, 满足
\[
 f(x)\le\lambda\quad\text{对几乎每个 }x\notin\bigcup_jQ_j,
 \qquad
 \lambda<\frac1{w(Q_j)}\int_{Q_j}fw\le C\lambda.
\]
证明本质上与 Lebesgue 测度的情形相同, 见定理~\ref{thm:ch2-calderon-zygmund-decomposition}. 有了这一点, 对 $w^{-1}$ 作 Calderón--Zygmund 分解, 即可像定理~\ref{thm:ch7-reverse-holder} 那样完成 \eqref{eq:ch7-dual-measure-reverse-holder} 的证明.

若 $w\in A_\infty$, 则由上述论证和命题~\ref{prop:ch7-ap-properties}, 存在 $p>1$, 使对任意 $q>p$, $w$ 以一致常数满足 $A_q$ 条件. 令 $q\to\infty$, 得
\begin{equation}
\label{eq:ch7-reverse-jensen}
 \frac1{|Q|}\int_Qw
 \le C\exp\left(\frac1{|Q|}\int_Q\log w\right).
\end{equation}
不等式 \eqref{eq:ch7-reverse-jensen} 称为反向 Jensen 不等式. 反之, 若 $w$ 对所有立方体 $Q$ 都满足 \eqref{eq:ch7-reverse-jensen}, 则 $w\in A_\infty$. 这一结论由 García-Cuerva 和 Rubio de Francia \MBUnresolvedReference{citation}{[6, p. 405]} 以及 S. V. Hruscev \MBUnresolvedReference{citation}{[A description of weights satisfying the $A_\infty$ condition of Muckenhoupt, Proc. Amer. Math. Soc. 90 (1984), 253--257]} 分别独立证明.

\subsubsection{奇异积分中 $A_p$ 条件的必要性}
\label{subsec:ch7-ap-necessity-singular-integrals}

$A_p$ 条件是 Hardy--Littlewood 极大函数在 $L^p(w)$, $1<p<\infty$, 上有界以及关于 $w$ 为弱型 $(1,1)$ 的充要条件. 对 Calderón--Zygmund 算子, 本章只证明了它是充分条件. 在如下意义下 $A_p$ 条件也是必要的: 若 $w$ 是 $\mathbb R^n$ 上的权, 且每个 Riesz 变换关于 $w$ 都是弱型 $(p,p)$ 的, $1<p<\infty$, 则 $w\in A_p$. 特别地, Hilbert 变换关于 $w$ 是弱型 $(p,p)$ 的当且仅当 $w\in A_p$. Hilbert 变换的结论见上引 Hunt, Muckenhoupt 和 Wheeden 的论文, 他们的论证可用于 $\mathbb R^n$ 上的 Riesz 变换. 此外 Stein \MBUnresolvedReference{citation}{[17, p. 210]} 证明: 若任一 Riesz 变换在 $L^p(w)$ 上有界, $1<p<\infty$, 则 $w\in A_p$.

\subsubsection{$A_p$ 权的因子分解}
\label{subsec:ch7-ap-factorization}

命题~\ref{prop:ch7-ap-properties} 的第 (3) 条允许用 $A_1$ 权构造 $A_p$ 权. 深刻的反命题也成立: $w\in A_p$ 当且仅当
\[
 w=w_0w_1^{1-p},
\]
其中 $w_0,w_1\in A_1$. 这称为因子分解定理, 最初由 Muckenhoupt 猜想, 后由 P. Jones 在 \MBUnresolvedReference{citation}{[Factorization of $A_p$ weights, Ann. of Math. 111 (1980), 511--530]} 中证明.

\hypertarget{chapter-07-page-163}{}
R. Coifman, P. Jones 和 J. L. Rubio de Francia 后来给出了一个很简单的证明, 见上引论文或 Stein \MBUnresolvedReference{citation}{[17, Chapter 5]}. J. L. Rubio de Francia 在上引论文中研究了更一般背景下的因子分解, 并证明了把算子因子分解与带权不等式联系起来的一般原则, 另见 \MBUnresolvedReference{citation}{[6, Chapter 6]}. 因子分解定理还可推广, 以包含关于反向 Hölder 不等式指数大小的信息. 见 D. Cruz-Uribe 和 C. J. Neugebauer 的论文 \MBUnresolvedReference{citation}{[The structure of the reverse Hölder classes, Trans. Amer. Math. Soc. 347 (1995), 2941--2960]}.

\subsubsection{$A_p$ 权与 BMO}
\label{subsec:ch7-ap-bmo}

设 $f$ 局部可积且 $\exp(f)\in A_2$. 给定立方体 $Q$, 记 $f_Q$ 为 $f$ 在 $Q$ 上的平均值. $A_2$ 条件给出
\[
 \left(\frac1{|Q|}\int_Q\exp(f)\right)
 \left(\frac1{|Q|}\int_Q\exp(-f)\right)\le C,
\]
等价地,
\[
 \left(\frac1{|Q|}\int_Q\exp(f-f_Q)\right)
 \left(\frac1{|Q|}\int_Q\exp(f_Q-f)\right)\le C.
\]
由 Jensen 不等式, 每个因子都至少为 $1$ 且至多为 $C$. 因而
\[
 \frac1{|Q|}\int_Q\exp(|f-f_Q|)\le2C,
\]
从而
\[
 \frac1{|Q|}\int_Q|f-f_Q|\le2C.
\]
故 $f\in\mathrm{BMO}$. 这证明了 $\exp(f)\in A_2$ 蕴含 $f\in\mathrm{BMO}$; 等价地, 若 $w\in A_2$, 则 $\log w\in\mathrm{BMO}$.

由 John--Nirenberg 不等式可证明: 若 $f\in\mathrm{BMO}$, 则对所有充分小的 $\lambda>0$, $\exp(\lambda f)\in A_2$, 比较推论~\ref{cor:ch6-bmo-exponential-integrability}. 换言之,
\[
 \mathrm{BMO}=\{\lambda\log w:\lambda\in\mathbb R,\ w\in A_2\}.
\]
事实上, 由命题~\ref{prop:ch7-ap-properties}, 这里的 $A_2$ 可换成任意 $A_p$, $1<p<\infty$.

\hypertarget{chapter-07-page-164}{}
\subsubsection{Coifman--Fefferman 不等式}
\label{subsec:ch7-coifman-fefferman}

Coifman 和 Fefferman 对定理~\ref{thm:ch7-cz-weighted-strong} 的原始证明依赖下述不等式. 对满足梯度条件 \eqref{eq:ch5-gradient-condition} 的奇异积分 $T$, 若 $w\in A_\infty$, 则对 $0<p<\infty$,
\begin{equation}
\label{eq:ch7-coifman-fefferman}
 \int_{\mathbb R^n}|T^*f|^pw
 \le C_p\int_{\mathbb R^n}|Mf|^pw.
\end{equation}
这里 $T^*$ 是与 $T$ 相伴的极大算子. 证明的核心是联系 $M$ 与 $T^*$ 的 good-$\lambda$ 不等式: 存在 $\delta>0$, 使对所有 $\gamma>0$,
\[
\begin{aligned}
 &w(\{x\in\mathbb R^n:|T^*f(x)|>2\lambda,\ Mf(x)<\gamma\lambda\})\\
 &\qquad\le C\gamma^\delta
 w(\{x\in\mathbb R^n:|T^*f(x)|>\lambda\}).
\end{aligned}
\]

R. Bagby 和 D. Kurtz 在 \MBUnresolvedReference{citation}{[A rearranged good $\lambda$ inequality, Trans. Amer. Math. Soc. 293 (1986), 71--81]} 中给出了 \eqref{eq:ch7-coifman-fefferman} 的另一证明, 并得到渐近尖锐的常数. 他们以重排不等式
\[
 (T^*f)_w^*(t)\le C(Mf)_w^*(t/2)+(Tf)_w^*(2t)
\]
替代 good-$\lambda$ 不等式, 其中 $(g)_w^*$ 表示 $g$ 关于测度 $w\,dx$ 的递减重排, 比较第~\ref{subsec:ch2-hardy-operator}~节. J. Alvarez 和 C. Perez 在 \MBUnresolvedReference{citation}{[Estimates with $A_\infty$ weights for various singular integral operators, Boll. Unione Mat. Ital. (7) 8-A (1994), 123--133]} 中把 \eqref{eq:ch7-coifman-fefferman} 推广到 Calderón--Zygmund 算子, 但把左端的 $T^*$ 换成 $T$. 他们证明了引理~\ref{lem:ch7-cz-sharp-maximal} 的更强形式: 对 $0<\delta<1$,
\[
 M^\#(|Tf|^\delta)(x)^{1/\delta}\le C_\delta Mf(x).
\]
于是 \eqref{eq:ch7-coifman-fefferman} 可由引理~\ref{lem:ch7-weighted-dyadic-sharp} 得到.

\subsubsection{强极大函数}
\label{subsec:ch7-strong-maximal-function}

在第~\ref{subsec:ch2-strong-maximal}~节中, 我们把强极大算子 $M_S$ 定义为函数在坐标轴平行矩形上的平均值上确界. 它满足加权范数不等式, 所需权类似于 $A_p$ 权. 强 $A_p$ 条件如下: 当 $1<p<\infty$ 时, $w\in A_p^*$ 是指对每个坐标轴平行矩形 $R$,
\[
 \left(\frac1{|R|}\int_Rw\right)
 \left(\frac1{|R|}\int_Rw^{1-p'}\right)^{p-1}\le C.
\]

\begin{Theorem}
\label{thm:ch7-strong-maximal-weighted}
设 $1<p<\infty$. $M_S$ 在 $L^p(w)$ 上有界当且仅当 $w\in A_p^*$.
\end{Theorem}

\hypertarget{chapter-07-page-165}{}
定理~\ref{thm:ch7-strong-maximal-weighted} 是 Hardy--Littlewood 极大函数相应结论的推论. 对 $1\le i\le n$, 定义只作用于第 $i$ 个变量的极大算子 $M_i$:
\[
 M_if(x_1,\ldots,x_n)
 =M(f(x_1,\ldots,x_{i-1},\cdot,x_{i+1},\ldots,x_n))(x_i).
\]
由 Fubini 定理, $M_Sf(x)\le(M_1\circ\cdots\circ M_n)f(x)$. 通过极限论证, 若 $w\in A_p^*$, 则对每个 $i$, $w(x_1,\ldots,x_{i-1},\cdot,x_{i+1},\ldots,x_n)$ 以一致常数满足一维 $A_p$ 条件. 此时定理~\ref{thm:ch7-strong-maximal-weighted} 立即由定理~\ref{thm:ch7-maximal-weighted-strong} 得到.

定义 $A_1^*=\{w:M_Sw\le Cw\text{ a.e.}\}$. $A_p^*$ 权可以用 $A_1^*$ 权分解.

\begin{Theorem}
\label{thm:ch7-strong-ap-factorization}
设 $1<p<\infty$. $w\in A_p^*$ 当且仅当存在 $w_0,w_1\in A_1^*$, 使
\[
 w=w_0w_1^{1-p}.
\]
\end{Theorem}

该定理首先由 J. L. Rubio de Francia 证明, 见第~\ref{subsec:ch7-references}~节所引论文. 然而定理~\ref{thm:ch7-a1-characterization} 的类似结论不成立; 当 $0<\delta<1$ 时, $(M_Sf)^\delta$ 未必属于 $A_1^*$. 这由 F. Soria 在 \MBUnresolvedReference{citation}{[A remark on $A_1$-weights for the strong maximal function, Proc. Amer. Math. Soc. 100 (1987), 46--48]} 中证明.

B. Jawerth 在 \MBUnresolvedReference{citation}{[Weighted inequalities for maximal operators: linearization, localization and factorization, Amer. J. Math. 108 (1986), 361--414]} 中把定理~\ref{thm:ch7-strong-maximal-weighted} 作为更一般结果的推论给出了另一证明. 设 $\mathcal B$ 是基, 即开集族. 给定权 $w$, 定义关于 $\mathcal B$ 和 $w$ 的带权极大函数
\[
 M_{\mathcal B,w}f(x)
 =\sup_{x\in B\in\mathcal B}\frac1{w(B)}\int_B|f|w,
\]
其中 $x\in\bigcup_{B\in\mathcal B}B$; 在其他点定义为 $0$. 当 $w=1$ 时简记为 $M_\mathcal B$. 若 $\mathcal B$ 是所有坐标轴平行矩形的集合, 则 $M_\mathcal B=M_S$. 若对每个 $B\in\mathcal B$,
\[
 \left(\frac1{|B|}\int_Bw\right)
 \left(\frac1{|B|}\int_Bw^{1-p'}\right)^{p-1}\le C,
\]
则称 $w$ 关于 $\mathcal B$ 满足 $A_p$ 条件, 记作 $w\in A_{p,\mathcal B}$.

\begin{Theorem}
\label{thm:ch7-basis-weighted-equivalence}
给定基 $\mathcal B$, 权 $w$ 和 $1<p<\infty$, 令 $\sigma=w^{1-p'}$. 下列条件等价:
\begin{enumerate}[label=(\arabic*)]
\item $M_\mathcal B$ 在 $L^p(w)$ 和 $L^{p'}(\sigma)$ 上都有界;
\item $w\in A_{p,\mathcal B}$, $M_{\mathcal B,w}$ 在 $L^{p'}(w)$ 上有界, 且 $M_{\mathcal B,\sigma}$ 在 $L^p(\sigma)$ 上有界.
\end{enumerate}
\end{Theorem}

\hypertarget{chapter-07-page-166}{}
$M_{\mathcal B,w}$ 和 $M_{\mathcal B,\sigma}$ 的有界性取决于基 $\mathcal B$ 的几何. 当 $\mathcal B$ 是矩形基时, 两个算子的有界性由 R. Fefferman 在 \MBUnresolvedReference{citation}{[Strong differentiation with respect to measures, Amer. J. Math. 103 (1981), 33--40]} 中以及 B. Jawerth 和 A. Torchinsky 在 \MBUnresolvedReference{citation}{[The strong maximal function with respect to measures, Studia Math. 80 (1984), 261--285]} 中证明. 对更一般的基, $M_{\mathcal B,w}$ 在 $L^{p'}(w)$ 上有界当且仅当该基满足一个与积分微分问题密切相关的覆盖性质. 详见上引 Jawerth 的论文.

\subsubsection{双权范数不等式}
\label{subsec:ch7-two-weights}

把本章结果推广到带不同权的 $L^p$ 空间是一个很困难的问题. 更确切地说, 给定算子 $T$, 要确定权对 $(u,v)$ 的充要条件, 使 $T$ 从 $L^p(v)$ 到 $L^p(u)$ 有界, 或满足弱型 $(p,p)$ 不等式
\[
 u(\{x\in\mathbb R^n:|Tf(x)|>\lambda\})
 \le\frac C{\lambda^p}\int_{\mathbb R^n}|f|^pv.
\]
对 Hardy--Littlewood 极大函数, 这些问题已有完整答案. 把 $A_p$ 条件推广如下: 当 $1<p<\infty$ 时, 若
\[
 \left(\frac1{|Q|}\int_Qu\right)
 \left(\frac1{|Q|}\int_Qv^{1-p'}\right)^{p-1}\le C,
\]
则称权对 $(u,v)\in A_p$; 若
\[
 Mu(x)\le Cv(x)\quad\text{对几乎每个 }x\in\mathbb R^n,
\]
则称 $(u,v)\in A_1$. 当 $u=v$ 时, 这正是第~\ref{sec:ch7-ap-condition}~节定义的 $A_p$ 条件. 对定理~\ref{thm:ch7-maximal-weighted-weak} 的证明稍作修改即可得到:

\begin{Theorem}
\label{thm:ch7-two-weight-maximal-weak}
给定 $1\le p<\infty$, 弱型不等式
\[
 u(\{x\in\mathbb R^n:Mf(x)>\lambda\})
 \le\frac C{\lambda^p}\int_{\mathbb R^n}|f|^pv
\]
成立当且仅当 $(u,v)\in A_p$.
\end{Theorem}

但是对强型 $(p,p)$ 不等式, $A_p$ 条件是必要而非充分的: 存在 $(u,v)\in A_p$ 使强型不等式不成立. 事实上, $(u,Mu)\in A_1\subset A_p$, 因而 $((Mu)^{1-p'},u^{1-p'})\in A_{p'}$. 这里使用了与命题~\ref{prop:ch7-ap-properties} 证明相同的论证: $(u,v)\in A_p$ 与 $(v^{1-p'},u^{1-p'})\in A_{p'}$ 等价. 若取 $u=|f|$, 该权对的强型 $(p',p')$ 不等式将蕴含 $M$ 在 $L^1$ 上有界, 与命题~\ref{prop:ch2-maximal-not-l1} 矛盾.

尽管如此, 强型 $(p,p)$ 不等式仍有一个刻画.

\hypertarget{chapter-07-page-167}{}
\begin{Theorem}
\label{thm:ch7-sawyer-testing}
给定权对 $(u,v)$ 和 $1<p<\infty$, 下列条件等价:
\begin{enumerate}[label=(\arabic*)]
\item $M$ 是从 $L^p(v)$ 到 $L^p(u)$ 的有界算子;
\item 对任意立方体 $Q$,
\begin{equation}
\label{eq:ch7-sawyer-testing}
 \int_QM(v^{1-p'}\chi_Q)^pu
 \le C\int_Qv^{1-p'}<\infty.
\end{equation}
\end{enumerate}
\end{Theorem}

不等式 \eqref{eq:ch7-sawyer-testing} 等价于 $M$ 在测试函数族 $v^{1-p'}\chi_Q$ 上有界. 定理~\ref{thm:ch7-sawyer-testing} 首先由 E. Sawyer 在 \MBUnresolvedReference{citation}{[A characterization of a two-weight norm inequality for maximal operators, Studia Math. 75 (1982), 1--11]} 中证明. D. Cruz-Uribe 在 \MBUnresolvedReference{citation}{[New proofs of two-weight norm inequalities for the maximal operator, Georgian Math. J. 7 (2000), 33--42]} 中给出了较简单的证明. 上一节所引 Jawerth 的论文给出了很不相同的证明. 当 $u=v$ 时, \eqref{eq:ch7-sawyer-testing} 等价于 $A_p$ 条件, 因而该结果也给出了定理~\ref{thm:ch7-maximal-weighted-strong} 的另一证明. 这由 R. Hunt, D. Kurtz 和 C. J. Neugebauer 在 \MBUnresolvedReference{citation}{[A note on the equivalence of $A_p$ and Sawyer's condition for equal weights, Conference on Harmonic Analysis in Honor of Antoni Zygmund, W. Beckner et al., eds., vol. 1, pp. 156--158, Wadsworth, Belmont, 1983]} 中证明.

对奇异积分算子, 除 Hilbert 变换外几乎没有已知结果. 当 $1<p<\infty$ 时, $A_p$ 条件对弱型 $(p,p)$ 不等式是必要但不充分的, 见 B. Muckenhoupt 和 R. Wheeden 的 \MBUnresolvedReference{citation}{[Two weight function norm inequalities for the Hardy--Littlewood maximal function and the Hilbert transform, Studia Math. 60 (1976), 279--294]}. 当 $p=1$ 时, Muckenhoupt 猜想 $A_1$ 条件是 Hilbert 变换弱型 $(1,1)$ 不等式的充要条件, 另见下一节. Hilbert 变换强型 $(2,2)$ 不等式已有充要条件; 更确切地说, 单位圆上的共轭函数 $\widetilde f$ 有这样的条件, 见第~\ref{subsec:ch3-conjugate-function-fourier-series}~节. 该结果归功于 M. Cotlar 和 C. Sadosky, 见 \MBUnresolvedReference{citation}{[On the Helson--Szegő theorem and a related class of modified Toeplitz kernels, Harmonic Analysis in Euclidean Spaces, G. Weiss and S. Wainger, eds., vol. 1, pp. 383--407, Proc. Sympos. Pure Math. 35, Amer. Math. Soc., Providence, 1979]}.

\subsubsection{双权不等式的进一步结果}
\label{subsec:ch7-more-two-weights}

应用带权不等式的一般技术可概括如下: 给定算子 $T$, 先找到正算子 $S$, 使 $T$ 从 $L^p(Su)$ 到 $L^p(u)$ 有界, 再由 $S$ 的有界性推出 $T$ 的有界性质. \MBUnresolvedReference{section}{第 8 章第 3 节}将使用这一技术.

\hypertarget{chapter-07-page-168}{}
为应用该技术, 需要找到这类双权不等式. 定理~\ref{thm:ch2-weighted-maximal} 就是一个例子: Hardy--Littlewood 极大函数从 $L^p(Mu)$ 到 $L^p(u)$ 有界. A. Córdoba 和 C. Fefferman 首次尝试把这一结果推广到其他算子, 见 \MBUnresolvedReference{citation}{[A weighted norm inequality for singular integrals, Studia Math. 57 (1976), 97--101]}. 他们证明, 若 $T$ 是满足梯度条件 \eqref{eq:ch5-gradient-condition} 的奇异积分, 则 $T$ 从 $L^p(M(u^s)^{1/s})$ 到 $L^p(u)$ 有界. 由定理~\ref{thm:ch7-a1-characterization}, $M(u^s)^{1/s}\in A_1\subset A_p$, 所以这一结果现在也由定理~\ref{thm:ch7-cz-weighted-strong} 得到.

下一定理给出更尖锐的结果.

\begin{Theorem}
\label{thm:ch7-iterated-maximal-singular-integral}
设 $T$ 是满足定理~\ref{thm:ch5-calderon-zygmund} 假设的奇异积分. 给定权 $u$, 当 $1<p<\infty$ 时,
\[
 \int_{\mathbb R^n}|Tf|^pu
 \le C_p\int_{\mathbb R^n}|f|^pM^{[p]+1}u,
\]
其中 $[p]$ 是 $p$ 的整数部分, $M^k=M\circ\cdots\circ M$ 是极大算子的第 $k$ 次迭代. 此不等式是尖锐的, 因为不能把 $[p]+1$ 换成 $[p]$.
\end{Theorem}

特别要注意, 奇异积分算子一般并不从 $L^p(Mu)$ 到 $L^p(u)$ 有界. 由于 $M(u^s)^{1/s}\in A_1$,
\[
 M^ku(x)\le M^k(M(u^s)^{1/s})(x)
 \le C_kM(u^s)(x)^{1/s},
\]
所以定理~\ref{thm:ch7-iterated-maximal-singular-integral} 强于 Córdoba 和 Fefferman 的结果. 该定理在 $1<p<2$ 时首先由 J. M. Wilson 在 \MBUnresolvedReference{citation}{[Weighted norm inequalities for the continuous square function, Trans. Amer. Math. Soc. 314 (1989), 661--692]} 中证明; C. Perez 在 \MBUnresolvedReference{citation}{[Weighted norm inequalities for singular integral operators, J. London Math. Soc. 49 (1994), 296--308]} 中证明了所有 $p>1$ 的情形. 在同一篇论文中 Perez 还证明了相应的弱型 $(1,1)$ 不等式:
\[
 u(\{x\in\mathbb R^n:|Tf(x)|>\lambda\})
 \le\frac C\lambda\int_{\mathbb R^n}|f|M^2u.
\]
该不等式能否把 $M^2$ 换成 $M$ 尚不清楚.
